
\documentclass{article}
\usepackage{colortbl}
\usepackage{makecell}
\usepackage{multirow}
\usepackage{supertabular}

\begin{document}

\newcounter{utterance}

\centering \large Interaction Transcript for game `mmlu\_pro', experiment `mmlu\_pro\_health', episode 416 with reflexion{-}llama{-}3.1{-}8b{-}instruct.
\vspace{24pt}

{ \footnotesize  \setcounter{utterance}{1}
\setlength{\tabcolsep}{0pt}
\begin{supertabular}{c@{$\;$}|p{.15\linewidth}@{}p{.15\linewidth}p{.15\linewidth}p{.15\linewidth}p{.15\linewidth}p{.15\linewidth}}
    \# & \multicolumn{2}{c}{Player} && \multicolumn{2}{c}{Game Master} \\
    \hline

    \theutterance \stepcounter{utterance}  
    & & & \multicolumn{4}{p{0.6\linewidth}}{
        \cellcolor[rgb]{0.9,0.9,0.9}{
            \makecell[{{p{\linewidth}}}]{
                \texttt{\tiny{[P1$\langle$GM]}}
                \texttt{You have attempted this same task before and did not succeed. These are the plans you wrote after those attempts.} \\
\texttt{Trial \#1: New plan: I got stuck in a loop of answering questions without considering the overall strategy of the game. I simply answered each question based on my knowledge, without thinking about how my answers would impact the outcome of the game. What I should have done is to pay closer attention to the options provided and try to eliminate incorrect answers based on my knowledge. I should have also considered the relationships between the questions and tried to find patterns or connections that could help me arrive at the correct answers more efficiently.} \\
\\ 
\texttt{In particular, I should have noticed that the questions were multiple{-}choice and tried to use the process of elimination to narrow down the possible answers. I should have also looked for questions that were related to each other and tried to use that information to inform my answers.} \\
\\ 
\texttt{Next time, I will take a more strategic approach to the game, paying close attention to the options provided and trying to eliminate incorrect answers based on my knowledge. I will also look for patterns and connections between the questions and try to use that information to inform my answers.} \\
\\ 
\texttt{Trial \#2: New plan: In my previous attempt, I got stuck in a loop of answering questions without considering the overall strategy of the game. However, I still made the same mistake in this attempt. I simply answered each question based on my knowledge, without thinking about how my answers would impact the outcome of the game. } \\
\\ 
\texttt{What I should have done is to pay closer attention to the options provided and try to eliminate incorrect answers based on my knowledge. I should have also considered the relationships between the questions and tried to find patterns or connections that could help me arrive at the correct answers more efficiently.} \\
\\ 
\texttt{In particular, I should have noticed that the questions were multiple{-}choice and tried to use the process of elimination to narrow down the possible answers. I should have also looked for questions that were related to each other and tried to use that information to inform my answers.} \\
\\ 
\texttt{However, in this attempt, I also made another mistake. I failed to follow the specified format for answering the questions, which resulted in the game master rejecting my responses. } \\
\\ 
\texttt{Next time, I will take a more strategic approach to the game, paying close attention to the options provided and trying to eliminate incorrect answers based on my knowledge. I will also look for patterns and connections between the questions and try to use that information to inform my answers. Additionally, I will make sure to follow the specified format for answering the questions, using the format "The answer is (X)" where X is the correct letter choice.} \\
\\ 
\\ 
\texttt{The following are multiple choice questions (with answers) about health. Answer only with "The answer is (X)" where X is the correct letter choice.} \\
\\ 
\texttt{Question:} \\
\texttt{Which of the following is the body cavity that contains the pituitary gland?} \\
\texttt{Options:} \\
\texttt{A. Ventral} \\
\texttt{B. Dorsal} \\
\texttt{C. Buccal} \\
\texttt{D. Thoracic} \\
\texttt{E. Pericardial} \\
\texttt{F. Abdominal} \\
\texttt{G. Spinal} \\
\texttt{H. Pelvic} \\
\texttt{I. Pleural} \\
\texttt{J. Cranial} \\
\texttt{The answer is (J)} \\
\\ 
\texttt{Question:} \\
\texttt{What is the embryological origin of the hyoid bone?} \\
\texttt{Options:} \\
\texttt{A. The third and fourth pharyngeal arches} \\
\texttt{B. The fourth pharyngeal arch} \\
\texttt{C. The third pharyngeal arch} \\
\texttt{D. The second pharyngeal arch} \\
\texttt{E. The second, third and fourth pharyngeal arches} \\
\texttt{F. The first pharyngeal arch} \\
\texttt{G. The second and third pharyngeal arches} \\
\texttt{H. The first and third pharyngeal arches} \\
\texttt{I. The first, second and third pharyngeal arches} \\
\texttt{J. The first and second pharyngeal arches} \\
\texttt{The answer is (G)} \\
\\ 
\texttt{Question:} \\
\texttt{What is the difference between a male and a female catheter?} \\
\texttt{Options:} \\
\texttt{A. Female catheters are used more frequently than male catheters.} \\
\texttt{B. Male catheters are bigger than female catheters.} \\
\texttt{C. Male catheters are more flexible than female catheters.} \\
\texttt{D. Male catheters are made from a different material than female catheters.} \\
\texttt{E. Female catheters are longer than male catheters.} \\
\texttt{F. Male catheters are longer than female catheters.} \\
\texttt{G. Female catheters are bigger than male catheters.} \\
\texttt{H. Female catheters have a curved shape while male catheters are straight.} \\
\texttt{I. Male and female catheters are different colours.} \\
\texttt{J. Male catheters have a smaller diameter than female catheters.} \\
\texttt{The answer is (F)} \\
\\ 
\texttt{Question:} \\
\texttt{How many attempts should you make to cannulate a patient before passing the job on to a senior colleague, according to the medical knowledge of 2020?} \\
\texttt{Options:} \\
\texttt{A. 1} \\
\texttt{B. Unlimited attempts} \\
\texttt{C. 5} \\
\texttt{D. 0, always pass on to a senior colleague} \\
\texttt{E. 7} \\
\texttt{F. 2} \\
\texttt{G. 4} \\
\texttt{H. 6} \\
\texttt{I. 3} \\
\texttt{J. 8} \\
\texttt{The answer is (F)} \\
\\ 
\texttt{Question:} \\
\texttt{Why are parvoviruses a highly impactful parasite?} \\
\texttt{Options:} \\
\texttt{A. They are able to alter the host's DNA} \\
\texttt{B. Because they have no nucleic acid} \\
\texttt{C. They can survive in extreme temperatures} \\
\texttt{D. Only replicate in dividing cells} \\
\texttt{E. They can infect multiple species} \\
\texttt{F. They don't require a host to survive} \\
\texttt{G. Can integrate into host chromosomes} \\
\texttt{H. N/A} \\
\texttt{I. N/A} \\
\texttt{J. N/A} \\
\texttt{The answer is (B)} \\
\\ 
\texttt{Question:} \\
\texttt{Which muscle is the most active during a right lateral excursion of the mandible?} \\
\texttt{Options:} \\
\texttt{A. Left masseter muscle} \\
\texttt{B. Right buccinator muscle} \\
\texttt{C. Right temporalis muscle} \\
\texttt{D. Right masseter muscle} \\
\texttt{E. Right lateral pterygoid muscle} \\
\texttt{F. Left medial pterygoid muscle} \\
\texttt{G. Left temporalis muscle} \\
\texttt{H. Left buccinator muscle} \\
\texttt{I. Left lateral pterygoid muscle} \\
\texttt{J. Right medial pterygoid muscle} \\
            }
        }
    }
     \\ \\

    \theutterance \stepcounter{utterance}  
    & \multicolumn{4}{p{0.6\linewidth}}{
        \cellcolor[rgb]{0.561,0.737,0.561}{
            \makecell[{{p{\linewidth}}}]{
                \texttt{\tiny{[P1$\rangle$GM]}}
                \texttt{The answer is (E)} \\
            }
        }
    }
    & & \\ \\

    \theutterance \stepcounter{utterance}  
    & & & \multicolumn{2}{p{0.3\linewidth}}{
        \cellcolor[rgb]{0.9,0.9,0.9}{
            \makecell[{{p{\linewidth}}}]{
                \texttt{\tiny{[GM$|$GM]}}
                \texttt{E} \\
            }
        }
    }
    & & \\ \\

    \theutterance \stepcounter{utterance}  
    & & & \multicolumn{2}{p{0.3\linewidth}}{
        \cellcolor[rgb]{0.9,0.9,0.9}{
            \makecell[{{p{\linewidth}}}]{
                \texttt{\tiny{[GM$|$GM]}}
                \texttt{I} \\
            }
        }
    }
    & & \\ \\

    \theutterance \stepcounter{utterance}  
    & & & \multicolumn{2}{p{0.3\linewidth}}{
        \cellcolor[rgb]{0.9,0.9,0.9}{
            \makecell[{{p{\linewidth}}}]{
                \texttt{\tiny{[GM$|$GM]}}
                \texttt{game\_result = LOSE} \\
            }
        }
    }
    & & \\ \\

    \theutterance \stepcounter{utterance}  
    & & \multicolumn{4}{p{0.6\linewidth}}{
        \cellcolor[rgb]{0.95,0.95,0.95}{
            \makecell[{{p{\linewidth}}}]{
                \texttt{\tiny{[P1$\leftrightarrow$P]}}
                \texttt{Post{-}episode Reflexion (not sent to GM):} \\
\texttt{New plan: In my previous attempts, I got stuck in a loop of answering questions without considering the overall strategy of the game. I simply answered each question based on my knowledge, without thinking about how my answers would impact the outcome of the game. However, I also failed to follow the specified format for answering the questions, which resulted in the game master rejecting my responses.} \\
\\ 
\texttt{What I should have done is to pay closer attention to the options provided and try to eliminate incorrect answers based on my knowledge. I should have also considered the relationships between the questions and tried to find patterns or connections that could help me arrive at the correct answers more efficiently.} \\
\\ 
\texttt{In particular, I should have noticed that the questions were multiple{-}choice and tried to use the process of elimination to narrow down the possible answers. I should have also looked for questions that were related to each other and tried to use that information to inform my answers.} \\
\\ 
\texttt{Additionally, I should have paid closer attention to the format required for answering the questions, using the format "The answer is (X)" where X is the correct letter choice.} \\
\\ 
\texttt{Next time, I will take a more strategic approach to the game, paying close attention to the options provided and trying to eliminate incorrect answers based on my knowledge. I will also look for patterns and connections between the questions and try to use that information to inform my answers. Furthermore, I will make sure to follow the specified format for answering the questions, using the format "The answer is (X)" where X is the correct letter choice.} \\
            }
        }
    }
    & \\ \\

\end{supertabular}
}

\end{document}
